% titel: Blatt 2, A2: Federpendel
% skript: Kap. Schwingungen, harmonischer Oszillator

% Federkraft, Ruhelage, Bewegungsgleichung, freie Schwingung.

Eine Masse an einer Feder schwingt um ihre Ruhelage mit der
Kreisfrequenz $\omega=\sqrt{k/m}$, solange die Auslenkungen das
Hooke'sche Gesetz erfuellen.

\begin{werkzeugkasten}[Werkzeuge]
\textbf{Hooke'sche Gesetz} (Kraft proportional zur Auslenkung gegen die
Ruhelage):
\begin{equation}
  F=-k\,(z-z_0),
\end{equation}
wobei $k$ die Federkonstante, $z_0$ die Ruhelage ist.

\textbf{Vertikale Feder:} Ruhelage ergibt sich aus dem Kräftegleichgewicht
(Gewicht gegen Federkraft)
\begin{equation}
  m g=k\,(z_0-z_{\text{unbelastet}})\;\Rightarrow\;
  z_0=\frac{m g}{k}\;.
\end{equation}

\textbf{Bewegungsgleichung} (2. Newton):
\begin{equation}
  m\ddot z=-k\,(z-z_0)\;.
\end{equation}
Mit $u=z-z_0$ wird daraus der harmonische Oszillator
\begin{equation}
  \ddot u+\omega^2 u=0,\qquad \omega=\sqrt{k/m},\quad
  T=2\pi/\omega\;.
\end{equation}

\textbf{Allgemeine Lösung} (frei wählbare $A,B$, festgelegt durch
Anfangsbedingungen $u(0)$ und $\dot u(0)$):
\begin{equation}
  u(t)=A\cos(\omega t)+B\sin(\omega t)=C\cos(\omega t+\varphi)\;.
\end{equation}
\end{werkzeugkasten}

\begin{werkzeugkasten}[Vorgehen]
\begin{itemize}
  \item Ruhelage: $F=0$ setzen bzw. Kräftegleichgewicht.
  \item Die DGL auf $u=z-z_0$ umschreiben, dann $A$, $B$ aus $z(0)$ und
        $\dot z(0)$ bestimmen.
  \item Ohne Anfangsgeschwindigkeit ($\dot u(0)=0$) bleibt nur der
        Kosinus-Term mit Amplitude $u(0)$.
\end{itemize}
\end{werkzeugkasten}

\begin{achtung}[Typische Fallen]
\begin{itemize}
  \item Die Ruhelage ist $z_0\ne0$; man rechnet sie sauber als konstanten
        Offset, bevor man die Oszillator-Formel ansetzt.
  \item $\omega$ ist \emph{nicht} $m/k$, sondern $\sqrt{k/m}$.
  \item Anfangsbedingungen inkenlich ansetzen: die Kugel wird \emph{nach
        unten} gedrückt; das Vorzeichen von $z(0)$ folgt aus der
        gewählten $z$-Achse.
\end{itemize}
\end{achtung}
